\section{Results}
\label{sec:results}

\subsection{Flexural Mode Frequencies Across Fill Volume}
\label{sec:freq_vs_volume}

Figure~\ref{fig:freq_vs_volume} presents the $n = 2$ and $n = 3$ flexural
mode frequencies across the physiological fill range
(\SIrange{50}{500}{\milli\litre}).  The $n = 2$ frequency follows a
distinctive U-shaped curve, with a minimum of \SI{13.5}{\hertz}
at \SI{222}{\milli\litre} and higher values at both extremes of the fill
range.

The physical mechanism behind this non-monotonic behaviour is a competition
between two opposing effects:
\begin{enumerate}
  \item \textbf{Geometric softening.}  As the bladder fills, the radius
    increases (equation~\ref{eq:radius}), which lowers the natural frequency
    through the $R^{-4}$ and $R^{-2}$ dependence in the bending and membrane
    stiffness terms (equations~\ref{eq:Kbend}--\ref{eq:Kmemb}).  The wall
    also thins (equation~\ref{eq:h_solve}), further reducing both bending
    rigidity ($D \propto h^3$) and membrane stiffness ($\propto h$).
  \item \textbf{Strain-stiffening.}  The elastic modulus of the detrusor
    muscle increases exponentially with fill volume
    (equation~\ref{eq:E_fill}), from \SI{29}{\kilo\pascal} at
    \SI{50}{\milli\litre} to \SI{145}{\kilo\pascal} at
    \SI{500}{\milli\litre}.  This five-fold increase eventually dominates
    the geometric effects.
\end{enumerate}
At low fill volumes the geometric effect dominates and the frequency
falls; at high fill volumes the strain-stiffening wins and the frequency
rises.  The crossover produces the minimum near \SI{222}{\milli\litre}---a
fill state corresponding to the early ``need to go'' sensation in most
individuals.

\begin{figure}[htbp]
  \centering
  \includegraphics[width=\textwidth]{../figures/fig_bladder_frequency_vs_volume.png}
  \caption{Bladder resonant properties across fill volume.  (a)~Flexural
  mode frequencies $f_2$ and $f_3$, with the ISO~2631 pelvic resonance band
  (\SIrange{4}{8}{\hertz}) shown for reference.  (b)~Geometry: equivalent
  sphere radius and wall thickness.  (c)~Wall elastic modulus showing
  exponential strain-stiffening.  (d)~Intravesical pressure following the
  cystometric curve.  The minimum $f_2 = \SI{13.5}{\hertz}$ occurs at
  \SI{222}{\milli\litre}, where geometric softening and strain-stiffening
  are in balance.}
  \label{fig:freq_vs_volume}
\end{figure}

Table~\ref{tab:frequencies} lists the computed frequencies at key fill states.
For fill states where the thin-shell approximation is valid
(\SIrange{150}{500}{\milli\litre}), the $n = 2$ mode spans
\SIrange{13.5}{17.1}{\hertz} and the $n = 3$ mode spans
\SIrange{22.4}{30.6}{\hertz}.  All predicted
frequencies lie above the ISO~2631 peak pelvic transmissibility band
(\SIrange{4}{8}{\hertz}), though below \SI{20}{\hertz} they remain within
the broader frequency range affected by whole-body vibration.

\begin{table}[htbp]
  \centering
  \caption{Flexural mode frequencies at key fill volumes.}
  \label{tab:frequencies}
  \begin{tabular}{rcccccc}
    \toprule
    Volume & $R$ & $h$ & $E$ & $f_2$ & $f_3$ & $f_4$ \\
    (mL) & (cm) & (mm) & (kPa) & (Hz) & (Hz) & (Hz) \\
    \midrule
     50 & 2.29 & 5.00 &  29.0 & 22.7 & 36.3 & 53.9 \\
    100 & 2.88 & 3.46 &  34.7 & 16.2 & 25.4 & 35.9 \\
    150 & 3.30 & 2.73 &  41.5 & 14.2 & 22.8 & 32.3 \\
    200 & 3.63 & 2.30 &  49.6 & 13.5 & 22.4 & 32.1 \\
    300 & 4.15 & 1.79 &  70.9 & 13.9 & 24.1 & 35.3 \\
    400 & 4.57 & 1.50 & 101.4 & 15.3 & 27.1 & 40.3 \\
    500 & 4.92 & 1.50 & 145.0 & 17.1 & 30.6 & 45.6 \\
    \bottomrule
  \end{tabular}
\end{table}


\subsection{Stiffness and Mass Decomposition}
\label{sec:decomposition}

To elucidate the physical mechanisms behind the U-shaped $f_2(V)$ curve, we
decompose the effective stiffness $K_\mathrm{total}$ and effective mass
$m_\mathrm{eff}$ into their constituent terms as functions of fill volume
(figure~\ref{fig:decomposition}).

At the frequency minimum ($V^* \approx \SI{222}{\milli\litre}$), the
stiffness is dominated by the pre-stress term:
$K_P = 58.7\%$, $K_\mathrm{memb} = 40.7\%$, and
$K_\mathrm{bend} = 0.6\%$.  Bending stiffness, which scales as $Dh^3/R^4$
and dominates thin-shell vibration in vacuo, is negligible for the
fluid-loaded bladder because the fluid added mass greatly exceeds the wall
mass.  The effective mass is similarly dominated by the fluid:
$m_\mathrm{fluid} = \rho_f R / n = 89.4\%$ of $m_\mathrm{eff}$, with the
wall contributing only $m_\mathrm{wall} = \rho_w h = 10.6\%$.

\paragraph{Analytical condition for the $f_2$ minimum.}
Since $\omega_n^2 = K_\mathrm{total} / m_\mathrm{eff}$, the condition
$\partial f_2 / \partial V = 0$ requires
\begin{equation}\label{eq:f2_min}
  \frac{\mathrm{d} \ln K_\mathrm{total}}{\mathrm{d}V}
  = \frac{\mathrm{d} \ln m_\mathrm{eff}}{\mathrm{d}V},
\end{equation}
i.e.\ the fractional growth rate of total stiffness must equal the
fractional growth rate of effective mass.  Figure~\ref{fig:decomposition}(d)
plots both log-rates; they cross at $V \approx \SI{222}{\milli\litre}$,
confirming the numerical minimum.  Below $V^*$, geometric softening causes
$K_\mathrm{total}$ to decrease faster (in fractional terms) than
$m_\mathrm{eff}$ increases, so $f_2$ falls.  Above $V^*$,
strain-stiffening and the rising intravesical pressure drive
$K_\mathrm{total}$ upward faster than the mass grows, so $f_2$ rises.

\begin{figure}[htbp]
  \centering
  \includegraphics[width=\textwidth]{figures/fig_stiffness_mass_decomposition.pdf}
  \caption{Stiffness and mass decomposition of the $n = 2$ mode.
  (a)~Stiffness components: bending ($K_\mathrm{bend}$, blue), membrane
  ($K_\mathrm{memb}$, green), pre-stress ($K_P$, red), and total (dashed).
  Pre-stress dominates above \SI{100}{\milli\litre}.
  (b)~Mass components: wall mass ($\rho_w h$, blue) and fluid added mass
  ($\rho_f R/n$, red).  The fluid contribution exceeds 85\% at all fill
  volumes.  (c)~Resulting $f_2$ curve with the minimum at
  $V^* = \SI{222}{\milli\litre}$, $f_2^\mathrm{min} = \SI{13.5}{\hertz}$.
  (d)~Logarithmic growth rates $\mathrm{d}(\ln K)/\mathrm{d}V$ and
  $\mathrm{d}(\ln m)/\mathrm{d}V$; their crossing at \SI{222}{\milli\litre}
  satisfies equation~\eqref{eq:f2_min}.}
  \label{fig:decomposition}
\end{figure}


\subsection{Sensitivity and Uncertainty Analysis}
\label{sec:sensitivity}

We assess the sensitivity of $f_2$ to five uncertain parameters using
one-at-a-time perturbation over their literature ranges
(table~\ref{tab:sensitivity}).  At the baseline volume
$V = \SI{222}{\milli\litre}$ (the frequency minimum), the wall
elastic modulus $E$ dominates all other parameters, with
$\Delta f_2 = \SI{14.9}{\hertz}$ (111\% of baseline) across the range
\SIrange{10}{500}{\kilo\pascal}.  The wall thickness $h$ is a distant
second at $\Delta f_2 = \SI{3.4}{\hertz}$ (25\%).  The Poisson ratio
changes $f_2$ by only \SI{0.14}{\hertz}, urine density by
\SI{0.06}{\hertz}, and the loss tangent not at all; these effects are
effectively negligible.

\begin{table}[htbp]
  \centering
  \caption{Sensitivity of $f_2$ to parameter uncertainty at
  $V = \SI{222}{\milli\litre}$ (baseline $f_2 = \SI{13.5}{\hertz}$).}
  \label{tab:sensitivity}
  \begin{tabular}{lccccc}
    \toprule
    Parameter & Range & $f_2^\mathrm{lo}$ & $f_2^\mathrm{hi}$ &
      $\Delta f_2$ & Fraction \\
    & & (Hz) & (Hz) & (Hz) & (\%) \\
    \midrule
    $E$ (wall modulus) & 10--500~kPa & 11.0 & 28.4 & 14.9 & 111 \\
    $h$ (wall thickness) & 2--6~mm & 13.3 & 16.8 & 3.4 & 25 \\
    $\nu$ (Poisson ratio) & 0.40--0.50 & 13.4 & 13.5 & 0.1 & 1 \\
    $\rho_f$ (urine density) & 1010--1030~kg/m$^3$ & 13.6 & 13.4 & 0.1 & $<$1 \\
    $\eta$ (loss tangent) & 0.30--0.50 & 13.5 & 13.5 & 0.0 & 0 \\
    \bottomrule
  \end{tabular}
\end{table}

Figure~\ref{fig:tornado} presents the tornado chart corresponding to
table~\ref{tab:sensitivity}.  The overwhelming dominance of $E$ reflects
the fact that both $K_\mathrm{memb}$ and $K_P$ (which together constitute
$>$99\% of the stiffness) are proportional to $E$ and to the intravesical
pressure, both of which are functions of the same strain-stiffening
mechanism.  The negligible sensitivity to $\eta$ confirms that damping
affects only the peak amplitude (via $Q$), not the resonant frequency.

\begin{figure}[htbp]
  \centering
  \includegraphics[width=0.85\textwidth]{figures/fig_tornado_sensitivity.pdf}
  \caption{Tornado chart showing the sensitivity of $f_2$ to each uncertain
  parameter at $V = \SI{222}{\milli\litre}$.  The wall elastic modulus $E$
  dominates, producing a \SI{14.9}{\hertz} swing across its literature
  range; wall thickness contributes \SI{3.4}{\hertz}, and all remaining
  parameters contribute $\leq$\SI{0.14}{\hertz}.}
  \label{fig:tornado}
\end{figure}

\paragraph{Robustness of the minimum.}
Figure~\ref{fig:min_shift} examines how the location $V^*$ and value
$f_2^\mathrm{min}$ of the frequency minimum shift under parameter
perturbation.  The minimum is remarkably robust to all parameters except
$E$: perturbing $h$ shifts $V^*$ from \SI{186}{\milli\litre} to
\SI{216}{\milli\litre} and $f_2^\mathrm{min}$ from \SI{13.2}{\hertz} to
\SI{16.8}{\hertz}, whereas $\nu$ and $\rho_f$ move the minimum by less than
\SI{3}{\milli\litre} and \SI{0.2}{\hertz}, and $\eta$ leaves it unchanged.
The elastic modulus, however, shifts the minimum substantially: at the upper
bound ($E = \SI{500}{\kilo\pascal}$), the minimum moves to
$V^* = \SI{462}{\milli\litre}$ with
$f_2^\mathrm{min} = \SI{21.2}{\hertz}$.  This shift occurs because a
stiffer wall delays the crossover between geometric softening and
strain-stiffening to higher volumes where the radius is larger.

The clinical implication is clear: individual variation in detrusor muscle
stiffness---the parameter with the widest literature range---is the
principal source of uncertainty in predicting both the minimum frequency
and the fill volume at which it occurs.

\begin{figure}[htbp]
  \centering
  \includegraphics[width=\textwidth]{figures/fig_minimum_shift.pdf}
  \caption{Robustness of the $f_2$ minimum to parameter uncertainty.
  (a)~Shift in $V^*$ (volume at minimum).
  (b)~Shift in $f_2^\mathrm{min}$ (frequency at minimum).
  Only $E$ produces a very large shift; $h$ causes a secondary but visible
  change, while $\nu$, $\rho_f$, and $\eta$ leave the minimum close to the
  baseline.}
  \label{fig:min_shift}
\end{figure}


\subsection{Coupling Pathways}
\label{sec:coupling_results}

Figure~\ref{fig:coupling} compares the airborne and mechanical coupling
pathways to the bladder at a representative fill of \SI{300}{\milli\litre}
($f_2 = \SI{13.9}{\hertz}$).  Three key observations emerge.

First, the airborne acoustic coupling scales as $(ka)^2$ and is vanishingly
small across the entire frequency range of interest: at \SI{13.9}{\hertz},
$(ka)^2 \approx 1.1 \times 10^{-4}$.  The bladder, like the abdominal
cavity, is deep in the Rayleigh regime at all physiologically relevant
frequencies.

Second, the mechanical coupling---the product of seat-to-pelvis
transmissibility and bladder modal amplification---peaks near the bladder
resonance frequency.  Even though pelvic transmissibility is attenuated
at \SI{13.9}{\hertz} ($T \approx 0.29$), the bladder modal amplification
($Q = 2.5$) partially compensates.

Third, the coupling ratio $\mathcal{R} \approx 6{,}400$ at $f_2$
(equation~\ref{eq:R_coupling}) confirms that mechanical vibration is the
dominant excitation pathway by nearly four orders of magnitude.  This is
quantitatively consistent with---though somewhat smaller than---the
$6.6 \times 10^4$ ratio found for the abdominal
cavity~\cite{Mace2025browntone}, the difference arising primarily from
the higher $f_2$ (and hence smaller $Q$ contribution from pelvic
transmissibility) of the bladder.

\begin{figure}[htbp]
  \centering
  \includegraphics[width=\textwidth]{../figures/fig_bladder_coupling.png}
  \caption{Coupling pathways to the bladder at \SI{300}{\milli\litre} fill.
  (a)~Airborne $(ka)^2$ coupling (blue) vs mechanical seat--pelvis--bladder
  coupling (red), with the $n{=}2$ resonance frequency marked.  (b)~Ratio of
  mechanical to airborne coupling (mechanical advantage), reaching
  $\sim$$10^4$--$10^6$ across the frequency range.  (c)~Components of the
  mechanical pathway: seat-to-pelvis transmissibility (green) and bladder
  modal amplification (magenta).}
  \label{fig:coupling}
\end{figure}


\subsection{Comparison with the Abdominal Cavity}
\label{sec:comparison}

Table~\ref{tab:comparison} compares key quantities between the bladder and
abdominal cavity models.  The bladder is smaller, softer, and more heavily
damped, yet the two organs share the same qualitative physics: fluid-dominated
breathing modes at kilohertz frequencies, low-frequency flexural modes
governed by the competition between shell stiffness and fluid added mass, and
a severe airborne coupling penalty in the Rayleigh regime.

\begin{table}[htbp]
  \centering
  \caption{Comparison of bladder and abdominal cavity models.}
  \label{tab:comparison}
  \begin{tabular}{lcc}
    \toprule
    Quantity & Bladder (300~mL) & Abdomen \\
    \midrule
    Equivalent radius & \SI{4.2}{\centi\metre} & \SI{15.7}{\centi\metre} \\
    Wall thickness & \SI{1.8}{\milli\metre} & \SI{10}{\milli\metre} \\
    Elastic modulus & \SI{71}{\kilo\pascal} & \SI{100}{\kilo\pascal} \\
    Poisson's ratio & 0.49 & 0.45 \\
    Loss tangent & 0.40 & 0.25 \\
    Quality factor $Q$ & 2.5 & 4.0 \\
    Breathing mode $f_0$ & $\sim$\SI{9500}{\hertz} & $\sim$\SI{2500}{\hertz} \\
    Flexural $f_2$ & \SI{13.9}{\hertz} & \SI{4.0}{\hertz} \\
    $ka$ at $f_2$ & 0.011 & 0.011 \\
    Coupling ratio $\mathcal{R}$ & $\sim$6,400 & $\sim$66,000 \\
    \bottomrule
  \end{tabular}
\end{table}

The principal quantitative differences are the higher flexural frequency of
the bladder ($f_2 = \SI{13.9}{\hertz}$ vs \SI{4.0}{\hertz}) and the smaller
coupling ratio ($\sim$6,400 vs $\sim$66,000).  The higher $f_2$ is a direct
consequence of the smaller radius: the stiffness-to-mass ratio scales
roughly as $1/R^2$, and the bladder is approximately four times smaller
than the abdomen.  The smaller coupling ratio reflects the fact that the
bladder's $f_2$ sits above the pelvic transmissibility peak, so the
mechanical pathway receives less amplification from the body's own
resonance structure.
